% Section 3.4 -- the instrument reference line, compressed. Diagnostics behind
% the precision term live in Appendix C. No literal numbers: macros only.

\subsection{The instrument reference line}
\label{sec:reference-line}

An extractor that recovers part of the explicit set displaces any comparison
built on its output, so the displacement has to be quantified. We do this with
a reference line, the human--machine gap the instrument produces on its own.

Two artifacts contribute and they are opposite in kind. Recall failure drops
premises a declaration really cites; precision failure admits premises it does
not. They do not cancel, because the corpora differ in the proof styles each
artifact is worst on, so both displace the comparison the same way and the line
is their sum. We therefore report it as separate terms rather than as one
averaged scalar.

The recall term is standardized jointly on tactic class and length bin, on the
step axis of Section~\ref{sec:length-axis}. It survives that control and
strengthens under it, from \recallTermClassOnly{} class-only to
\recallTermJoint{} joint, with a bootstrap interval of (\recallTermCIlow{},
\recallTermCIhigh{}) and a band of (\recallTermBandLow{},
\recallTermBandHigh{}) across the precedence orderings of the class partition.
The correction runs against our own interest, since a larger instrument
artifact raises the bar our headline result has to clear.

The precision term is the weak point. By a criterion registered before the run
it is uninformative rather than clean. We therefore carry it as an open item
rather than folding it into the line, and Appendix~\ref{app:reference} gives
the diagnostics behind that decision.

Two limits on what this line licenses. Per-class efficiency is close to flat
(Section~\ref{sec:retention}), which is a stronger neutrality statement than
the raw per-class miss rates suggest, but it is measured on human proofs,
and nothing here establishes comparable recovery on machine proofs, which
Section~\ref{sec:limitations}, L1, states as the paper's load-bearing
assumption. And the line matters unevenly: the coverage result of
Section~\ref{sec:coverage} sits far outside the regime these artifacts operate
in, while the secondary characterizations sit inside it.
